%
%Не забыть:
%--------------------------------------
%Вставить колонтитулы, поменять название на титульнике



%--------------------------------------

\documentclass[a4paper, 12pt]{article} 

%--------------------------------------
%Russian-specific packages
%--------------------------------------
%\usepackage[warn]{mathtext}
\usepackage[T2A]{fontenc}
\usepackage[utf8]{inputenc}
\usepackage[english,russian]{babel}
\usepackage[intlimits]{amsmath}
\usepackage{esint}
%--------------------------------------
%Hyphenation rules
%--------------------------------------
\usepackage{hyphenat}
\hyphenation{ма-те-ма-ти-ка вос-ста-нав-ли-вать}
%--------------------------------------
%Packages
%--------------------------------------
\usepackage{amsmath}
\usepackage{amssymb}
\usepackage{amsfonts}
\usepackage{amsthm}
\usepackage{latexsym}
\usepackage{mathtools}
\usepackage{etoolbox}%Булевые операторы
\usepackage{extsizes}%Выставление произвольного шрифта в \documentclass
\usepackage{geometry}%Разметка листа
\usepackage{indentfirst}
\usepackage{wrapfig}%Создание обтекаемых текстом объектов
\usepackage{fancyhdr}%Создание колонтитулов
\usepackage{setspace}%Настройка интерлиньяжа
\usepackage{lastpage}%Вывод номера последней страницы в документе, \lastpage
\usepackage{soul}%Изменение параметров начертания
\usepackage{hyperref}%Две строчки с настройкой гиперссылок внутри получаеммого
\usepackage[usenames,dvipsnames,svgnames,table,rgb]{xcolor}% pdf-документа
\usepackage{multicol}%Позволяет писать текст в несколько колонок
\usepackage{cite}%Работа с библиографией
\usepackage{subfigure}% Человеческая вставка нескольких картинок
\usepackage{tikz}%Рисование рисунков
\usepackage{float}% Возможность ставить H в положениях картинки
% Для картинок Моти
\usepackage{misccorr}
\usepackage{lscape}
\usepackage{cmap}
\usepackage{empheq}
\newcommand*\widefbox[1]{\fbox{\hspace{2em}#1\hspace{2em}}}



\usepackage{graphicx,xcolor}
\graphicspath{{Pictures/}}
\DeclareGraphicsExtensions{.pdf,.png,.jpg}

%----------------------------------------
%Список окружений
%----------------------------------------
\newenvironment {theor}[2]
{\smallskip \par \textbf{#1.} \textit{#2}  \par $\blacktriangleleft$}
{\flushright{$\blacktriangleright$} \medskip \par} %лемма/теорема с доказательством
\newenvironment {proofn}
{\par $\blacktriangleleft$}
{$\blacktriangleright$ \par} %доказательство
%----------------------------------------
%Список команд
%----------------------------------------
\newcommand{\grad}
{\mathop{\mathrm{grad}}\nolimits\,} %градиент

\newcommand{\diver}
{\mathop{\mathrm{div}}\nolimits\,} %дивергенция

\newcommand{\rot}
{\ensuremath{\mathrm{rot}}\,}

\newcommand{\Def}[1]
{\underline{\textbf{#1}}} %определение

\newcommand{\RN}[1]
{\MakeUppercase{\romannumeral #1}} %римские цифры

\newcommand {\theornp}[2]
{\textbf{#1.} \textit{ #2} \par} %Написание леммы/теоремы без доказательства

\newcommand{\qrq}
{\ensuremath{\quad \Rightarrow \quad}} %Человеческий знак следствия

\newcommand{\qlrq}
{\ensuremath{\quad \Leftrightarrow \quad}} %Человеческий знак равносильности

\renewcommand{\phi}{\varphi} %Нормальный знак фи

\newcommand{\me}
{\ensuremath{\mathbb{E}}}

\newcommand{\md}
{\ensuremath{\mathbb{D}}}

\newcommand{\bra}[1]
{\ensuremath{\left\langle#1\right|}}

\newcommand{\cat}[1]
{\ensuremath{\left|#1\right\rangle}}



%\renewcommand{\vec}{\overline}




%----------------------------------------
%Разметка листа
%----------------------------------------
\geometry{top = 3cm}
\geometry{bottom = 2cm}
\geometry{left = 1.5cm}
\geometry{right = 1.5cm}
%----------------------------------------
%Колонтитулы
%----------------------------------------
\pagestyle{fancy}%Создание колонтитулов
\fancyhead{}
%\fancyfoot{}
%\fancyhead[R]{\textsc{Уравнения Лондонов. Кинетическая индуктивность сверхпроводников.}}%Вставить колонтитул сюда
%----------------------------------------
%Интерлиньяж (расстояния между строчками)
%----------------------------------------
%\onehalfspacing -- интерлиньяж 1.5
%\doublespacing -- интерлиньяж 2
%----------------------------------------
%Настройка гиперссылок
%----------------------------------------
\hypersetup{				% Гиперссылки
	unicode=true,           % русские буквы в раздела PDF
	pdftitle={Заголовок},   % Заголовок
	pdfauthor={Автор},      % Автор
	pdfsubject={Тема},      % Тема
	pdfcreator={Создатель}, % Создатель
	pdfproducer={Производитель}, % Производитель
	pdfkeywords={keyword1} {key2} {key3}, % Ключевые слова
	colorlinks=true,       	% false: ссылки в рамках; true: цветные ссылки
	linkcolor=blue,          % внутренние ссылки
	citecolor=blue,        % на библиографию
	filecolor=magenta,      % на файлы
	urlcolor=red           % на URL
}
%----------------------------------------
%Работа с библиографией (как бич)
%----------------------------------------
\renewcommand{\refname}{Список литературы}%Изменение названия списка литературы для article
%\renewcommand{\bibname}{Список литературы}%Изменение названия списка литературы для book и report
%----------------------------------------
\begin{document}
	\begin{titlepage}
		\begin{center}
			$$$$
			$$$$
			$$$$
			$$$$
			{\Large{НАЦИОНАЛЬНЫЙ ИССЛЕДОВАТЕЛЬСКИЙ УНИВЕРСИТЕТ}}\\
			\vspace{0.1cm}
			{\Large{ВЫСШАЯ ШКОЛА ЭКОНОМИКИ}}\\
			\vspace{0.25cm}
			{\large{Факультет физики}}\\
			\vspace{5.5cm}
			{\Huge\textbf{{Домашнее задание}}}\\%Общее название
			\vspace{1cm}
			{\LARGE{Квантовая механика, неделя 15}}\\%Точное название
			\vspace{2cm}
			{Задание выполнил студент 2 курса}\\
			{Захаров Сергей Дмитриевич}
			\vfill
			\includegraphics[width = 0.2\textwidth]{HSElogo}\\
			\vfill
			Москва\\
			2020
		\end{center}
	\end{titlepage}

\section*{Задача 1}

Находить значение скалярного произведения будем с помощью рассмотрения $\hat{\mathbf{S}}^2$:

\begin{equation}
	\hat{\mathbf{S}}^2 = \hat{\mathbf{s}}_1^2 + \hat{\mathbf{s}}_2^2 + 2 \cdot  \hat{\mathbf{s}}_1 \cdot \hat{\mathbf{s}}_2 \qrq \hat{\mathbf{s}}_1 \cdot \hat{\mathbf{s}}_2 = \frac{1}{2} \left(\hat{\mathbf{S}}^2 - \hat{\mathbf{s}}_1^2 - \hat{\mathbf{s}}_2^2\right)
	\label{eq:t1_result}
\end{equation}

Для $\hat{\mathbf{s}}^2$ мы можем посчитать собственное значение как $s \cdot (s + 1)$, где $s$ --- собственное значение оператора $\hat{\mathbf{s}}$. Зная, что спины равны $1/2$  мы сразу же получаем, что собственные числа $\hat{\mathbf{s}}_1^2$ и $\hat{\mathbf{s}}_2^2$ равны $3/4$. 

Теперь необходимо определить собственное число оператора $\hat{\mathbf{S}}^2$. Для этого сперва найдем собственное число оператора $\hat{\mathbf{S}} = \hat{\mathbf{s}}_1 + \hat{\mathbf{s}}_2$. Спины могут быть либо сонаправлены (тогда они сложатся и собственное число будет равно 1), либо противоположно направлены (тогда собственное число будет равно 0). На основании этого делаем вывод, что собственные числа $\hat{\mathbf{S}}^2$ по уже указанной формуле равны либо $0$, либо $2$.

На основании этого, с помощью формулы (\ref{eq:t1_result}) получаем:

\begin{equation*}
	\boxed{
	\left[ \begin{array}{lr}
	\hat{\mathbf{s}}_1 \cdot \hat{\mathbf{s}}_2 = -\dfrac{3}{4}\\
	\hat{\mathbf{s}}_1 \cdot \hat{\mathbf{s}}_2 = \dfrac{1}{4}\\
	\end{array}
	\right.}
\end{equation*}

\section*{Задача 2}
\subsection*{а)}

Если мы будем рассматривать спиноры относительно оси $\mathbf{n}$, то, очевидным образом, они будут следующими:

\begin{equation}
	\alpha_{1\mathbf{n}} = 
	\begin{pmatrix}
		1\\
		0
	\end{pmatrix}
	\quad
	\alpha_{1\mathbf{n}} = 
	\begin{pmatrix}
		0\\
		1
	\end{pmatrix}
\end{equation} 

При этом, спиноры относительно оси $\mathbf{z}$ и $\mathbf{n}$ связаны следующим образом:

\begin{equation}
	\alpha_{\mathbf{n}} = \hat{U} \alpha_\mathbf{z}
	\label{eq:main_2}
\end{equation}

где в силу условия мы можем утверждать, что $\hat{U} = \hat{U}_z \hat{U}_y$. В таком случае:

\begin{align}
	\hat{U} = \hat{U}_y \hat{U}_z = 
	\begin{pmatrix}
		\cos\theta/2 & \sin\theta/2 \\
		-\sin\theta/2 & \cos\theta/2
	\end{pmatrix}
	\cdot
	\begin{pmatrix}
		\exp(i\phi/2) & 0 \\
		0 & \exp(-i\phi/2)
	\end{pmatrix}
	= \\
	=
	\begin{pmatrix}
		\exp(i\phi/2) \cos\theta/2 & \exp(-i \phi/2) \sin\theta/2 \\
		-\exp(i\phi/2) \sin\theta/2 & \exp(-i\phi/2) \cos\theta/2
	\end{pmatrix}
\end{align}

Теперь необходимо найти обратную матрицу к этой. В силу простоты этих вычислений (обращение матрицы 2x2), приведу сразу ответ:

\begin{equation}
	\hat{U}^{-1} = 
	\begin{pmatrix}
		\exp(-i\phi/2) \cos\theta/2 & -\exp(-i\phi/2) \sin\theta/2 \\
		\exp(i\phi/2) \sin\theta/2 & \exp(i\phi/2) \cos\theta/2
	\end{pmatrix}
\end{equation}

Тогда по формуле $\ref{eq:main_2}$ с учетом ее домножения слева на $\hat{U}^{-1}$:

\begin{empheq}[box=\widefbox]{align}
	\alpha_{1\mathbf{z}} &= \hat{U}^{-1} \alpha_{1\mathbf{n}} = 
	\begin{pmatrix}
		\exp(-i\phi/2) \cos\theta/2 \\
		\exp(i\phi/2) \sin\theta/2
	\end{pmatrix}\\
	\alpha_{2\mathbf{z}} &= \hat{U}^{-1} \alpha_{2\mathbf{n}} = 
	\begin{pmatrix}
	-\exp(-i\phi/2) \sin\theta/2 \\
	\exp(i\phi/2) \cos\theta/2
	\end{pmatrix}
\end{empheq}

\subsection*{б)}

Для начала сперва запишем выражения для вектора $\mathbf{n}$ через углы $\theta$ и $\phi$:

\begin{equation}
	\mathbf{n} = 
	\begin{pmatrix}
		\sin\theta \cos\phi \\
		\sin\theta \sin\phi \\
		\cos\theta
	\end{pmatrix}
	\label{eq:n}
\end{equation}

С учетом выражения для оператора спина через матрицы Паули мы в таком случае можем сказать, что матрица оператора проекции спина будет следующей:

\begin{equation}
	\sin\theta \cos\phi \cdot 
	\begin{pmatrix}
		0 & 1 \\
		1 & 0 
	\end{pmatrix}
	+ \sin\theta \sin\phi \cdot 
	\begin{pmatrix}
		0 & -i \\
		i & 0
	\end{pmatrix}
	+ \cos\theta \cdot
	\begin{pmatrix}
		1 & 0\\
		0 & -1
	\end{pmatrix}
	=
	\begin{pmatrix}
		\cos\theta & -\exp(-i \phi)\sin\theta \\
		\exp(i \phi)\sin\theta & -\cos\theta
	\end{pmatrix}
\end{equation}

Для того, чтобы найти интересующие нас спиноры, достаточно найти собственные векторы этой матрицы (в силу того, что проекции спина в проекции на необходимую нам ось составляют 1/2 и $-1/2$). Задача в целом тривиальная, поэтому сразу ответ:

\begin{empheq}[box=\widefbox]{align}
	\alpha_{1\mathbf{z}} &=
	\begin{pmatrix}
		\exp(-i \phi) \cdot \cos\theta/2 \\
		\cos\theta/2
	\end{pmatrix}
	= \exp(-i\phi/2) \cdot 
	\begin{pmatrix}
		\exp(-i\phi/2) \cos\theta/2 \\
		\exp(i\phi/2) \sin\theta/2
	\end{pmatrix}\\
	\alpha_{2\mathbf{z}} &= 
	\begin{pmatrix}
		-\exp(- i \phi) \cdot \sin\theta/2 \\
		\cos\theta/2
	\end{pmatrix}
	= \exp(-i\phi/2) \cdot
	\begin{pmatrix}
		-\exp(-i\phi/2) \sin\theta/2 \\
		\exp(i\phi/2) \cos\theta/2
	\end{pmatrix}
\end{empheq}


\end{document}